\documentclass[11pt,a4paper]{article}

% ---------------------------------------------------------------------------
% Packages
% ---------------------------------------------------------------------------
\usepackage[utf8]{inputenc}
\usepackage[T1]{fontenc}
\usepackage{graphicx} % Required for inserting images
\usepackage[margin=2.5cm]{geometry}
\usepackage{booktabs}
\usepackage{longtable}
\usepackage{array}
\usepackage{seqsplit}
\usepackage{xcolor}
\usepackage{listings}
\usepackage{hyperref}
\usepackage{enumitem}
\usepackage{parskip}

% ---------------------------------------------------------------------------
% Listings style (used for curl commands, JSON payloads, HCL, properties...)
% ---------------------------------------------------------------------------
\definecolor{codebg}{rgb}{0.96,0.96,0.96}
\definecolor{codecomment}{rgb}{0.45,0.45,0.45}
\definecolor{codestring}{rgb}{0.6,0.0,0.0}
\definecolor{codekeyword}{rgb}{0.0,0.0,0.6}

\lstdefinestyle{code}{
    backgroundcolor=\color{codebg},
    commentstyle=\color{codecomment}\itshape,
    keywordstyle=\color{codekeyword}\bfseries,
    stringstyle=\color{codestring},
    basicstyle=\ttfamily\footnotesize,
    breakatwhitespace=false,
    breaklines=true,
    captionpos=b,
    keepspaces=true,
    numbers=left,
    numbersep=6pt,
    numberstyle=\tiny\color{codecomment},
    showspaces=false,
    showstringspaces=false,
    showtabs=false,
    tabsize=2,
    frame=single,
    rulecolor=\color{black!20},
    xleftmargin=1em,
    xrightmargin=1em
}
\lstset{style=code}

\hypersetup{
    colorlinks=true,
    linkcolor=blue!50!black,
    urlcolor=blue!50!black,
    citecolor=blue!50!black,
    pdftitle={VPPaaS - Enterprise Integration - Sprint 2 Report},
    pdfauthor={Afonso Carvalho (IST 116482), Eduardo Silva (IST 106929)}
}

% ---------------------------------------------------------------------------
\title{
    \textbf{Virtual Power Plant as a Service (VPPaaS)} \\[0.4em]
    \Large Enterprise Integration --- Sprint 2 Technical Report \\[0.2em]
    \large Business Process Orchestration with Camunda~8 and Kong API Gateway
}
\author{
    Afonso Carvalho \textit{(IST 116482)} \\
    Eduardo Silva \textit{(IST 106929)}
}
\date{June 2026}

\begin{document}

\maketitle

\begin{abstract}
\noindent
This report documents the Sprint~2 deliverable of the VPPaaS Enterprise
Integration project. While Sprint~1 delivered the data plane (microservices,
shared RDS database, Kafka backbone), Sprint~2 \emph{orchestrates} it into the
eight business processes of the project brief using \textbf{Camunda~8} as the
process engine and \textbf{Kong API Gateway} as the single entry point to every
microservice. It covers (i) how each of the eight flows was modelled in BPMN
and wired to the microservices through Kong; (ii) their end-to-end execution
against the live AWS deployment, including the defects found and corrected
during testing; and (iii) the documentation of the source code, Terraform
infrastructure-as-code, installation procedures and parametrizations of the
integration layer.
\end{abstract}

\tableofcontents
\newpage

% ===========================================================================
\section{Introduction}
% ===========================================================================

The VPPaaS platform aggregates Distributed Energy Resources (DER) --- battery
storage systems (BESS), solar inverters (PV) and EV chargers --- and exposes
them as a single virtual power plant that a Utility Operator can use to balance
supply and demand across grid cells. Sprint~1 delivered the data plane ---
the domain microservices, the shared AWS RDS database and the Kafka event
backbone, all provisioned through Terraform --- documented in
\texttt{Report.md}/\texttt{Report.pdf}.

Sprint~2 closes the loop required by the project brief: \emph{``orchestrate the
business processes that coordinate these microservices, expose them through an
API Gateway, and demonstrate that the resulting flows execute correctly from
end to end''}. Concretely, the brief requires:

\begin{enumerate}[label=(\arabic*)]
    \item Modelling and deploying, through \textbf{Camunda} (orchestration
          engine) and \textbf{Kong} (API Gateway), the eight business
          processes that the platform must support: \emph{Prosumer
          Management}, \emph{Utility Operator Management}, \emph{Asset Link
          Management}, \emph{Telemetry Ingestion}, \emph{Flexibility
          Emission}, \emph{Grid Balancing Recommendation}, \emph{Energy
          Analytics} and \emph{Flexibility Forecasting (AI)};
    \item Documenting \textbf{end-to-end tests} of the full execution of these
          business processes against the deployed system;
    \item Documenting the \textbf{source code, Terraform scripts, BPMN models,
          installation procedures and parametrizations} that make the
          integration layer reproducible.
\end{enumerate}

Sections~\ref{sec:architecture}--\ref{sec:docs} below address these three
points in order.

% ===========================================================================
\section{Business Process Architecture}
\label{sec:architecture}
% ===========================================================================

\subsection{Integration platform: Camunda~8 and Kong API Gateway}

Two new infrastructure components were added on top of the Sprint~1 data plane:

\begin{itemize}
    \item \textbf{Camunda~8 (Zeebe + Operate + Tasklist), EC2 \texttt{t3.large}}
          --- the orchestration engine that executes the BPMN process
          definitions, exposes the \texttt{/v2/process-instances} and
          \texttt{/v2/deployments} REST APIs, renders generated forms in
          Tasklist and gives full execution visibility (active instances,
          completed/incident instances, variables) in Operate.
          Provisioned by \texttt{Camunda-Terraform/EC2CamundaEngine.tf}, an
          \texttt{aws\_instance} of type \texttt{t3.large} with an
          all-open security group (ports 8080/8081/8082), running the
          official Camunda~8.8.9 Docker Compose stack via the \texttt{deploy.sh}
          user-data script. Login: \texttt{demo}/\texttt{demo}.
    \item \textbf{Kong API Gateway, EC2 \texttt{t3.small}} --- the single
          public entry point for \emph{every} microservice. Provisioned by
          \texttt{KongTerraform/EC2Kong.tf} and configured at boot by
          \texttt{setup-kongGateway.sh}, which registers one Kong
          \emph{service} and one \emph{route} per microservice (see
          Table~\ref{tab:kong-routes}). Every BPMN service task that needs to
          reach a microservice calls Kong --- never the microservice's EC2
          address directly --- which is exactly the integration pattern the
          brief asks for (``expose the business processes through an API
          Gateway'').
    \item \textbf{Konga, EC2 \texttt{t3.small}} --- a web administration UI
          for Kong (\texttt{KongaTerraform/EC2Konga.tf}), used during
          development to inspect/validate the services and routes that
          \texttt{setup-kongGateway.sh} registers.
\end{itemize}

\renewcommand{\arraystretch}{1.25}
\begin{table}[h!]
\centering
\footnotesize
\caption{Kong services and routes registered by \texttt{setup-kongGateway.sh} (one per microservice)}
\label{tab:kong-routes}
\begin{tabular}{@{}p{0.33\textwidth}p{0.30\textwidth}p{0.24\textwidth}@{}}
\toprule
\textbf{Kong service name} & \textbf{Route path} (\texttt{strip\_path=true}) & \textbf{Upstream microservice} \\
\midrule
\texttt{\seqsplit{prosumer-service}} & \texttt{/prosumer} & Prosumer MS \\
\texttt{\seqsplit{utility-operator-service}} & \texttt{\seqsplit{/utility-operator}} & UtilityOperator MS \\
\texttt{\seqsplit{assetlink-service}} & \texttt{/assetlink} & AssetLink MS \\
\texttt{\seqsplit{telemetry-service}} & \texttt{/telemetry} & Telemetry MS \\
\texttt{\seqsplit{flexibility-event-service}} & \texttt{\seqsplit{/flexibility-event}} & FlexibilityEvent MS \\
\texttt{\seqsplit{grid-balancing-service}} & \texttt{\seqsplit{/grid-balancing}} & GridBalancing MS \\
\texttt{\seqsplit{energy-analytics-service}} & \texttt{\seqsplit{/energy-analytics}} & EnergyAnalytics MS \\
\texttt{\seqsplit{artificial-intelligence-service}} & \texttt{\seqsplit{/artificial-intelligence}} & ArtificialIntelligence MS \\
\bottomrule
\end{tabular}
\renewcommand{\arraystretch}{1}
\end{table}

\noindent
With \texttt{strip\_path=true}, a request such as
\texttt{GET http://<kong-dns>:8000/prosumer/Prosumer} arrives at the Prosumer
microservice as \texttt{GET /Prosumer}. All the BPMN flows described below use
this single Kong base URL --- in the current deployment
\url{http://ec2-100-54-67-134.compute-1.amazonaws.com:8000} --- as the target
of their HTTP Connector service tasks, so the orchestration layer never needs
to know individual microservice EC2 addresses.

\subsection{Modelling pattern: Initiator/Executor collaborations}

Two of the eight flows (\emph{Prosumer Management} and \emph{Asset Link
Management}/\emph{Telemetry Ingestion}) were modelled as \textbf{BPMN
collaborations} made up of several independently-deployable processes that
communicate through asynchronous \emph{messages}, following a classic
``Initiator $\rightarrow$ Executor'' pattern:

\begin{enumerate}
    \item The \emph{Initiator} process starts from a \texttt{none start event},
          collects input through a Tasklist user-task form (e.g. the data of
          the Prosumer to create), and then \textbf{publishes a message} that
          starts a second, independent process instance (the \emph{Executor}).
    \item The \emph{Executor} process is triggered by a \textbf{message start
          event}, runs the actual business logic --- promise/possibility
          checks, the REST call to the microservice through Kong, and the
          declaration step --- and finally publishes a confirmation message
          back to the Initiator.
    \item The Initiator resumes on a matching \textbf{message catch event} and
          completes with an acceptance user task.
\end{enumerate}

This mirrors the request/promise/declare/accept commitment lifecycle described
in the project brief's business-architecture section, and demonstrates
asynchronous, message-based collaboration between independently deployed
process definitions --- not just a single linear flow. The remaining flows
(Utility Operator Management, Telemetry/AI/Analytics/Flexibility/Balancing)
are modelled as single linear processes, since their business logic does not
require cross-participant negotiation.

A structural detail that is easy to overlook --- and that produced one of the
defects documented in Section~\ref{sec:bugs} --- is that \textbf{every variable
the Executor needs (in particular \texttt{camundaBaseUrl}, used by Tasklist
form-generation service tasks) must be explicitly forwarded in the
\texttt{variables} object of the message that starts it}; Camunda does not
automatically share process variables across independent process instances.

\subsection{Mapping the eight required business processes to BPMN models}

Table~\ref{tab:flows} maps each of the eight business processes required by
the brief to its BPMN file, Camunda \texttt{processDefinitionId}(s) and the
microservices it orchestrates (always reached through Kong, per
Table~\ref{tab:kong-routes}).

\renewcommand{\arraystretch}{1.35}
\begin{table}[h!]
\centering
\footnotesize
\caption{The eight required business processes and their BPMN realisation}
\label{tab:flows}
\begin{tabular}{@{}p{0.205\textwidth}p{0.205\textwidth}p{0.21\textwidth}p{0.255\textwidth}@{}}
\toprule
\textbf{Business process} & \textbf{BPMN file} & \textbf{processDefinitionId(s)} & \textbf{Microservices orchestrated} \\
\midrule
1. Prosumer Management & \texttt{\seqsplit{ProsumerManagement.bpmn}} & \texttt{\seqsplit{ProsumerMngInitiator}} + Executor process & Prosumer MS \\
2. Utility Operator Management & \texttt{\seqsplit{UtilityOperatorManagement.bpmn}} & \texttt{\seqsplit{UtilityOperatorManagement}} & UtilityOperator MS \\
3. Asset Link Management & \texttt{\seqsplit{AssetLinkManagement.bpmn}} & \texttt{\seqsplit{ProsumerForAssetLink}} $\to$ \texttt{\seqsplit{AssetLinkManagement}} & Prosumer MS, AssetLink MS \\
4. Telemetry Ingestion & \texttt{\seqsplit{AssetLinkManagement.bpmn}} (3rd participant) & \texttt{\seqsplit{TelemetryManagement}} & AssetLink MS, Telemetry MS \\
5. Flexibility Emission & \texttt{\seqsplit{FlexibilityEmission.bpmn}} & \texttt{\seqsplit{FlexibilityEmissionProcess}} & Telemetry MS, FlexibilityEvent MS \\
6. Grid Balancing Recommendation & \texttt{\seqsplit{GridBalancing.bpmn}} & \texttt{\seqsplit{GridBalancingProcess}} & UtilityOperator MS, Telemetry MS, AssetLink MS, Prosumer MS, GridBalancing MS \\
7. Energy Analytics & \texttt{\seqsplit{EnergyAnalytics.bpmn}} & \texttt{\seqsplit{EnergyAnalyticsProcess}} & Telemetry MS, Prosumer MS, UtilityOperator MS, EnergyAnalytics MS \\
8. Flexibility Forecasting (AI) & \texttt{\seqsplit{AIForecasting.bpmn}} & \texttt{\seqsplit{AIForecastingProcess}} & FlexibilityEvent MS, ArtificialIntelligence MS \\
\bottomrule
\end{tabular}
\end{table}
\renewcommand{\arraystretch}{1}

\noindent
Note that \emph{Asset Link Management} and \emph{Telemetry Ingestion} are
modelled inside the \textbf{same} BPMN file
(\texttt{AssetLinkManagement.bpmn}) as a three-participant collaboration
(\texttt{ProsumerForAssetLink} $\rightarrow$ \texttt{AssetLinkManagement}
$\rightarrow$ \texttt{TelemetryManagement}); this reflects the real dependency
described in the brief, where activating an asset link is what triggers the
creation of the dynamic Kafka consumer for that link's telemetry topic.

\subsection{Detailed description of each flow}

\subsubsection*{1 --- Prosumer Management}
\textbf{Collaboration} \texttt{ProsumerForAssetLink-style Initiator/Executor}.
The Initiator process \texttt{ProsumerMngInitiator} starts from a none start
event and presents the user task \emph{``Decide the Data for Prosumer Creation
order''} (generated form fields: \texttt{Name}, \texttt{Fiscal Number},
\texttt{Location}). The service task \emph{``Request Prosumer Creation
order''} then publishes a message carrying \texttt{name}, \texttt{location},
\texttt{fiscalnumber} and \texttt{camundaBaseUrl}, starting the Executor
process. The Executor runs \emph{``Verify if execute is possible''} (a
checkbox user task: \emph{``Is it possible to promise the request?''}),
\emph{``Promise Prosumer Creation order''}, the HTTP Connector service task
\emph{``Create Prosumer''} (\texttt{POST
http://<kong-dns>:8000/prosumer/Prosumer}), \emph{``Declare Prosumer Creation
order''} and finally \emph{``Accept Prosumer Creation order''}, after which
both processes reach their end events. This is the flow that was driven to
completion in the live environment; see Section~\ref{sec:e2e-prosumer}.

\subsubsection*{2 --- Utility Operator Management}
\textbf{Single linear process} \texttt{UtilityOperatorManagement}, with no
required start variables. It chains seven user/service tasks, each
\texttt{assignee=demo}: \emph{Request} $\rightarrow$ \emph{Verify if
... is possible} $\rightarrow$ \emph{Promise} $\rightarrow$ \emph{Create
Utility Operator} (real REST call, \texttt{POST
http://<kong-dns>:8000/utility-operator/UtilityOperator}) $\rightarrow$
\emph{Declare} $\rightarrow$ \emph{Check} $\rightarrow$ \emph{Accept}, each
gated by an \emph{``Is it possible/acceptable?''} checkbox form. Input data
collected on the first task: \texttt{Name}, \texttt{Location}.

\subsubsection*{3 --- Asset Link Management}
\textbf{Three-participant collaboration} inside
\texttt{AssetLinkManagement.bpmn}: \texttt{ProsumerForAssetLink} (Initiator,
none start event) $\rightarrow$ \texttt{AssetLinkManagement} (message start,
\texttt{messageRef = Message\_2r1j62e}) $\rightarrow$
\texttt{TelemetryManagement} (message start, \texttt{messageRef =
Message\_2mp8beu}). The Initiator's user task \emph{``Decide the data to
AssetLink association order''} (form \texttt{FormChoiceDataForAssetLink};
fields \texttt{UtilityOperatorID}, \texttt{prosumerID}) feeds a message that
starts \texttt{AssetLinkManagement}, which runs \emph{``Verify if execute
product is possible''} (form \texttt{Form\_Promise\_AssetLinkMng}, checkbox
\texttt{promise}) and the REST call that creates the AssetLink record (via
Kong, \texttt{/assetlink/AssetLink}); on completion it both replies to the
Initiator and starts \texttt{TelemetryManagement}. The Initiator finally
collects the acceptance user tasks \emph{``Check AssetLink association
order''} / \emph{``Check Telemetry Consumer Order''} (form
\texttt{Form\_Accept\_AssetLinkMng}, checkbox \texttt{accept}).

\subsubsection*{4 --- Telemetry Ingestion}
\textbf{Third participant of the same collaboration}, process
\texttt{TelemetryManagement}. Triggered by the message published by
\texttt{AssetLinkManagement} once an asset link is created, it runs its own
\emph{``Verify if execute product is possible''} step (same form/checkbox
pattern, \texttt{Form\_Promise\_AssetLinkMng}/\texttt{promise}) and the HTTP
Connector service task that calls \texttt{POST
http://<kong-dns>:8000/telemetry/Telemetry/Consume} to start the dynamic
Kafka consumer for the new \texttt{\{assetLinkId\}-\{utilityOperatorId\}}
topic --- the integration point between the orchestration layer and the
Telemetry MS's \texttt{DynamicTopicConsumer} described in Sprint~1
(\texttt{Report.md}, \S3.5). It then replies with a confirmation message
consumed back in the Initiator.

\subsubsection*{5 --- Flexibility Emission}
\textbf{Single linear process} \texttt{FlexibilityEmissionProcess}. The user
task \emph{``Select Asset for Flexibility Emission''} collects \texttt{Asset
ID} (e.g. \texttt{BATT-001}/\texttt{SOLAR-001}) and \texttt{Grid Cell ID}
(e.g. \texttt{LISBON-DT}); automatic service tasks then query the latest
telemetry for that asset through Kong (\texttt{/telemetry/Telemetry/asset/...})
and call \texttt{POST .../flexibility-event/Flexibility/Emit}, exercising the
rule-based emission logic of the FlexibilityEvent MS (D2 in
\texttt{Report.md}).

\subsubsection*{6 --- Grid Balancing Recommendation}
\textbf{Single linear process} \texttt{GridBalancingProcess}. The user task
\emph{``Select Grid Zone for Balancing''} collects \texttt{Grid Cell ID}
(\texttt{LISBON-DT}); the process then automatically aggregates data from the
UtilityOperator, Telemetry, AssetLink and Prosumer microservices (through
Kong) and calls \texttt{POST .../grid-balancing/GridBalancing/Emit},
reproducing the data-acquisition chain documented as D3 in \texttt{Report.md}.

\subsubsection*{7 --- Energy Analytics}
\textbf{Single linear process} \texttt{EnergyAnalyticsProcess}. The user task
\emph{``Select Zone for Energy Analytics''} collects \texttt{Grid Cell ID};
the process aggregates telemetry, prosumer and operator data through Kong and
calls \texttt{POST .../energy-analytics/EnergyAnalytics/Publish}, corresponding
to the D4 flow in \texttt{Report.md}.

\subsubsection*{8 --- Flexibility Forecasting (AI)}
\textbf{Single linear process} \texttt{AIForecastingProcess}. The user task
\emph{``Select Zone for AI Flexibility Forecast''} collects \texttt{Grid Cell
ID} (\texttt{LISBON-DT}); the automatic chain \emph{``Get Flexibility Event
History''} $\rightarrow$ \emph{``AI Analysis via Ollama (ArtificialIntelligence
MS)''} $\rightarrow$ \emph{``Emit AI Flexibility Event (AI\_DISCHARGE)''}
reproduces the D5 flow: it queries historical \texttt{FlexibilityEvent} rows,
sends them to the \texttt{ArtificialIntelligence} MS (Ollama / \texttt{llama3.2})
through Kong, and republishes the AI-enriched event with
\texttt{offerType = AI\_DISCHARGE}. As this performs a real LLM inference call,
its execution time is noticeably longer than the other flows.

% ===========================================================================
\section{End-to-end Testing}
\label{sec:testing}
% ===========================================================================

\subsection{Test environment and methodology}

All tests were executed against the live AWS deployment produced by
\texttt{DeploymentAutomation-ubuntu.sh} / \texttt{-macOS.sh} (Sprint~1
infrastructure plus the Camunda, Kong and Konga instances added for Sprint~2):

\begin{itemize}
    \item \textbf{Camunda~8}: \url{http://ec2-54-159-105-82.compute-1.amazonaws.com:8080}
          (Operate at \texttt{/operate}, login \texttt{demo}/\texttt{demo});
    \item \textbf{Kong Gateway}: \url{http://ec2-100-54-67-134.compute-1.amazonaws.com:8000};
    \item \textbf{Konga}: \url{http://ec2-98-92-246-7.compute-1.amazonaws.com:1337}.
\end{itemize}

\noindent
\textbf{Method.} Every flow with a \texttt{none start event} is started with a
\texttt{POST /v2/process-instances} call against the Camunda REST API (see
Listing~\ref{lst:start-process} for the generic template); the execution is
then driven step by step in \textbf{Operate}, completing the generated
Tasklist forms with realistic seed data (the \texttt{LISBON-DT} grid cell, the
\texttt{ArcoCegoLisbon} operator, the seed prosumer \texttt{client1} and assets
\texttt{BATT-001}/\texttt{SOLAR-001} created in Sprint~1, plus the new
\texttt{Jo\~ao Silva} prosumer created during this sprint's tests). Completion
is verified by querying the corresponding microservice \textbf{through Kong}
and checking that a new row was persisted in the shared RDS database. The full
script of commands and form values for all eight flows is kept in
\texttt{BPMN/EndToEnd\_Tests.md}, written specifically to support
the live defence demonstration.

\begin{lstlisting}[language=bash, caption={Generic template for starting any none-start-event process}, label={lst:start-process}]
curl -X POST "http://<camunda-dns>:8080/v2/process-instances" \
  -H "Content-Type: application/json" \
  -d '{"processDefinitionId": "<PROCESS_ID>", "variables": { ... }}'
\end{lstlisting}

\subsection{Full end-to-end execution: Prosumer Management}
\label{sec:e2e-prosumer}

The \emph{Prosumer Management} flow was driven to completion end to end ---
from the Camunda REST API, through every user task in Operate/Tasklist, through
Kong, into the Prosumer microservice, and finally into a persisted row in the
shared RDS MySQL database. This is the flow chosen to demonstrate, in the
defence, that the orchestration layer genuinely calls the deployed
microservices (rather than stopping at a mocked or partially-wired model).

\textbf{1. Start the Initiator instance:}
\begin{lstlisting}[language=bash]
curl -X POST "http://ec2-54-159-105-82.compute-1.amazonaws.com:8080/v2/process-instances" \
  -H "Content-Type: application/json" \
  -d '{"processDefinitionId": "ProsumerMngInitiator",
       "variables": {"camundaBaseUrl": "http://ec2-54-159-105-82.compute-1.amazonaws.com:8080"}}'
\end{lstlisting}

\textbf{2. Drive the flow in Operate/Tasklist} (Section~\ref{sec:architecture}
describes each step in detail): fill \emph{``Decide the Data for Prosumer
Creation order''} with \texttt{Name = Jo\~ao Silva}, \texttt{Fiscal Number =
123456789}, \texttt{Location = Lisboa}; the Executor process starts
automatically; complete \emph{``Verify if execute is possible''} by checking
\emph{``Is it possible to promise the request?''} = yes; the chain
\emph{Promise} $\rightarrow$ \emph{Create Prosumer} $\rightarrow$
\emph{Declare} $\rightarrow$ \emph{Accept} runs automatically and both process
instances reach their end events.

\textbf{3. Verify persistence through Kong:}
\begin{lstlisting}[language=bash, caption={Verification call and observed result --- a new row was created}]
$ curl http://ec2-100-54-67-134.compute-1.amazonaws.com:8000/prosumer/Prosumer
[
  {"FiscalNumber": 123456, "id": 1, "location": "Lisbon",  "name": "client1"},
  {"FiscalNumber": 123456789, "id": 2, "location": "Lisboa", "name": "Joao Silva"}
]
\end{lstlisting}

\noindent
Record \texttt{id = 2} is the row created live by the orchestrated flow,
confirming that the chain \textbf{Camunda $\to$ Kong $\to$ Prosumer
microservice $\to$ RDS MySQL} works end to end. (An earlier run produced a row
with \texttt{FiscalNumber: null}; that defect, and how it was corrected, is
documented as Defect~3 in Section~\ref{sec:bugs}, and the run shown above is
the corrected, fresh execution.)

\subsection{Defects identified and corrected}
\label{sec:bugs}

Driving the Prosumer flow end to end --- and then proactively re-checking the
remaining seven flows for the same classes of problem --- surfaced three
systemic defects in the BPMN models. None of them were visible from a static
read of the diagrams; all three only manifested at runtime, which is precisely
why end-to-end testing (rather than design-time validation alone) is required
by the brief. All three were fixed \emph{at the BPMN-definition level} and the
affected processes were \textbf{redeployed and re-run from a clean instance}
--- the user's explicit preference over patching a running incident through
Operate's variable panel, since a definition-level bug will simply recur on
the next instantiation if only the symptom is patched.

\begin{small}
\begin{longtable}{@{}p{0.20\textwidth}p{0.37\textwidth}p{0.37\textwidth}@{}}
\toprule
\textbf{Defect} & \textbf{Symptom \& root cause} & \textbf{Fix applied} \\
\midrule
\endhead
\textbf{1. \texttt{camundaBaseUrl} not propagated between collaborating processes}
&
Incident raised in Operate: \texttt{\seqsplit{jakarta.validation.ValidationException}: Property: url: ... must not be blank}. Root cause: the Initiator never forwarded its \texttt{camundaBaseUrl} variable inside the message body that starts the Executor, so the Executor's Tasklist form-generation service task received a blank URL. This is a \emph{structural} collaboration bug, not a one-off. &
Added \texttt{"camundaBaseUrl": camundaBaseUrl} as the first key of the \texttt{variables} object published by the message-sending service tasks: in \texttt{ProsumerManagement.bpmn} (``Request Prosumer Creation order'', Initiator~$\to$~Executor) and, proactively, in \emph{both} message-publication points of \texttt{AssetLinkManagement.bpmn} (\texttt{ProsumerForAssetLink}~$\to$~\texttt{AssetLinkManagement} and~$\to$~\texttt{TelemetryManagement}). \\
\addlinespace
\textbf{2. Empty Kong-address placeholder in service-task URLs}
&
Service tasks failed with a malformed target \texttt{\seqsplit{http://:8000/prosumer/Prosumer}} (empty host). Root cause: an unsubstituted \texttt{KONG\_ADDRESS} placeholder left over from the BPMN authoring/templating step --- present not just in Prosumer Management but systemically across the project. &
Found \textbf{21 occurrences across all 7 BPMN files} (1 in \texttt{ProsumerManagement}, fixed individually first; 20 more in \texttt{AIForecasting}, \texttt{AssetLinkManagement}, \texttt{EnergyAnalytics}, \texttt{UtilityOperatorManagement}, \texttt{GridBalancing}, \texttt{FlexibilityEmission}). Corrected with one global substitution (Listing~\ref{lst:sed-fix}) replacing the empty host with the live Kong DNS, verified afterwards with a zero-result grep for the broken pattern. \\
\addlinespace
\textbf{3. Variable name case-mismatch dropped the fiscal number}
&
The created Prosumer record persisted with \texttt{"FiscalNumber": null}. Root cause: the ``Create Prosumer'' service task referenced the FEEL variable \texttt{fiscalNumber} (camelCase), while the actual process variable produced by the form is \texttt{fiscalnumber} (lowercase) --- so the expression resolved to \texttt{null}. &
Corrected the request-body expression in \texttt{ProsumerManagement.bpmn} from \texttt{"FiscalNumber": fiscalNumber} to \texttt{"FiscalNumber": fiscalnumber}; redeployed and re-ran a fresh instance, which produced the correctly-populated record \texttt{id = 2} shown in Listing~\ref{lst:e2e-verify-result} (FiscalNumber~=~123456789). \\
\bottomrule
\caption{Defects found while end-to-end testing the BPMN flows, their root cause and the corrective action taken at the process-definition level.}
\label{tab:bugs}
\end{longtable}
\end{small}

\begin{lstlisting}[language=bash, caption={Single global substitution that corrected the empty-Kong-address defect (Defect 2) in 20 locations across 6 files in one pass}, label={lst:sed-fix}]
cd BPMN && sed -i \
  's|http://:8000|http://ec2-100-54-67-134.compute-1.amazonaws.com:8000|g' \
  AIForecasting.bpmn AssetLinkManagement.bpmn EnergyAnalytics.bpmn \
  UtilityOperatorManagement.bpmn GridBalancing.bpmn FlexibilityEmission.bpmn
\end{lstlisting}

\noindent
After every edit to a \texttt{.bpmn} file, the corrected definition must be
\textbf{redeployed} before being tested again (Listing~\ref{lst:redeploy}); a
running instance keeps executing the version of the definition it was started
with, so an in-place edit has no effect on instances already in flight.

\begin{lstlisting}[language=bash, caption={Redeploying a corrected BPMN definition to Camunda}, label={lst:redeploy}]
curl -L -X POST "http://ec2-54-159-105-82.compute-1.amazonaws.com:8080/v2/deployments" \
  -H "Content-Type: multipart/form-data" -H "Accept: application/json" \
  -F "resources=@./BPMN/<FILE_NAME>.bpmn"
\end{lstlisting}

\begin{lstlisting}[caption={Final verification result referenced from Section~\ref{sec:e2e-prosumer} (Defect 3 corrected: FiscalNumber persists)}, label={lst:e2e-verify-result}]
{"FiscalNumber": 123456789, "id": 2, "location": "Lisboa", "name": "Joao Silva"}
\end{lstlisting}

\subsection{Test procedures for the remaining seven flows}

To make the remaining flows reproducible for the live defence, every flow was
documented with its exact start command, its Tasklist step sequence (with
concrete seed-data values to type into each form) and its Kong verification
call, in \texttt{BPMN/EndToEnd\_Tests.md}. The flows are grouped by
modelling complexity:

\begin{itemize}
    \item \textbf{Group 1 --- single-process flows} (\emph{Utility Operator
          Management}, \emph{AI Forecasting}, \emph{Energy Analytics},
          \emph{Flexibility Emission}, \emph{Grid Balancing
          Recommendation}): no collaboration/messages, no
          \texttt{camundaBaseUrl} requirement --- only affected by Defect~2
          (broken Kong address), already corrected in all six files.
    \item \textbf{Group 2 --- \emph{Asset Link Management}/\emph{Telemetry
          Ingestion}}: equivalent in complexity to \emph{Prosumer
          Management} --- a three-participant message collaboration
          (\texttt{ProsumerForAssetLink} $\to$ \texttt{AssetLinkManagement}
          $\to$ \texttt{TelemetryManagement}), where Operate shows three
          separate process instances that must be alternately advanced; the
          same two corrections (Defects~1 and~2) were applied here too.
\end{itemize}

\noindent
For \emph{Group~1}, each flow follows the same shape: \texttt{POST
/v2/process-instances} with the bare \texttt{processDefinitionId} (no start
variables required), one user task that selects the test subject (a
\texttt{Grid Cell ID} of \texttt{LISBON-DT}, an \texttt{Asset ID} of
\texttt{BATT-001}/\texttt{SOLAR-001}, etc.), then an automatic chain of
service tasks that aggregates data from the relevant microservices through
Kong and posts the result to the owning microservice's ``emit/publish''
endpoint. \emph{Group~2} follows the Initiator/Executor pattern described in
Section~\ref{sec:architecture}, seeded with \texttt{UtilityOperatorID = 1}
(\texttt{ArcoCegoLisbon}) and \texttt{ProsumerID = 1} or \texttt{2} (the
\texttt{Jo\~ao Silva} record created in Section~\ref{sec:e2e-prosumer}).

% ===========================================================================
\section{Source Code, Infrastructure and Parametrization Documentation}
\label{sec:docs}
% ===========================================================================

This section documents the artefacts that are \emph{new} in Sprint~2 ---
the BPMN process definitions and the Camunda/Kong/Konga infrastructure that
orchestrates and exposes the Sprint~1 microservices. The microservice source
code, RDS/Kafka Terraform, deployment automation and Kafka parametrization
were already documented in full in \texttt{Report.md} (Sprint~1 deliverable);
that documentation is referenced rather than duplicated, and is summarised in
Table~\ref{tab:sprint1-recap} for completeness.

\subsection{Repository layout of the orchestration layer}

\begin{lstlisting}[caption={Layout of the new orchestration-layer artefacts inside the project repository}]
EI-Project/
 +-- BPMN/
 |    +-- ProsumerManagement.bpmn        # Flow 1 (Initiator/Executor collab.)
 |    +-- UtilityOperatorManagement.bpmn # Flow 2 (linear)
 |    +-- AssetLinkManagement.bpmn       # Flows 3 & 4 (3-participant collab.)
 |    +-- FlexibilityEmission.bpmn       # Flow 5 (linear)
 |    +-- GridBalancing.bpmn             # Flow 6 (linear)
 |    +-- EnergyAnalytics.bpmn           # Flow 7 (linear)
 |    +-- AIForecasting.bpmn             # Flow 8 (linear)
 |    +-- EndToEnd_Tests.md     # Step-by-step E2E test script (all 8 flows)
 +-- Camunda-Terraform/
 |    +-- EC2CamundaEngine.tf            # t3.large EC2 + security group
 |    +-- deploy.sh                      # user-data: boots Camunda 8.8.9 stack
 +-- KongTerraform/
 |    +-- EC2Kong.tf                     # t3.small EC2 + security group
 |    +-- deploy.sh                      # user-data: installs/starts Kong
 |    +-- setup-kongGateway.sh           # registers 8 services + 8 routes
 +-- KongaTerraform/
      +-- EC2Konga.tf                    # t3.small EC2 running the Konga admin UI
\end{lstlisting}

\subsection{Camunda~8 engine --- Terraform}

\texttt{Camunda-Terraform/EC2CamundaEngine.tf} provisions a single
\texttt{t3.large} \texttt{aws\_instance} (Camunda~8 needs more memory than the
\texttt{t3.small} tier used for the microservices, since it runs Zeebe,
Operate, Tasklist and Elasticsearch together) with an all-open security group
(ports~0--65535, \texttt{0.0.0.0/0} --- adequate for a coursework
demonstration, as already noted for the microservice security groups in
Sprint~1) and a \texttt{user\_data\_replace\_on\_change = true} script that
pulls and starts the official Camunda~8.8.9 Docker Compose distribution,
exposing Operate/Tasklist/the REST API on ports 8080/8081/8082.

\begin{lstlisting}[language=bash, caption={Core resource definition of the Camunda EC2 instance (\texttt{EC2CamundaEngine.tf})}]
resource "aws_instance" "exampleInstallCamundaEngine" {
  ami                    = "ami-07ff62358b87c7116"
  instance_type          = "t3.large"
  vpc_security_group_ids = [aws_security_group.instance.id]
  key_name               = "vockey"
  user_data              = "${file("deploy.sh")}"
  user_data_replace_on_change = true
}
\end{lstlisting}

\subsection{Kong API Gateway --- Terraform and service registration}

\texttt{KongTerraform/EC2Kong.tf} provisions the Kong instance (\texttt{t3.small},
same all-open security-group pattern) and exposes its public DNS as a
Terraform output (\texttt{address}), which the deployment automation captures
and forwards to \texttt{setup-kongGateway.sh}.

\begin{lstlisting}[language=bash, caption={Core resource definition of the Kong EC2 instance (\texttt{EC2Kong.tf})}]
resource "aws_instance" "exampleInstallKong" {
  ami                    = "ami-07ff62358b87c7116"
  instance_type          = "t3.small"
  vpc_security_group_ids = [aws_security_group.instance.id]
  key_name               = "vockey"
  user_data              = "${file("deploy.sh")}"
}

output "address" {
  value       = aws_instance.exampleInstallKong.public_dns
  description = "Connect to the KONG at this endpoint"
}
\end{lstlisting}

\noindent
\texttt{setup-kongGateway.sh} is the parametrization script that turns a bare
Kong instance into the platform's API Gateway: it receives the Kong DNS plus
the DNS of every microservice as positional parameters, polls
\texttt{http://<kong-dns>:8001/status} until the Admin API answers (up to 60
attempts, 10\,s apart), and then calls a \texttt{register\_service} helper
once per microservice that (a) \texttt{POST}s to \texttt{/services/} with the
service's upstream URL and (b) \texttt{POST}s to
\texttt{/services/\{name\}/routes} with \texttt{strip\_path=true} and the
route path shown in Table~\ref{tab:kong-routes}.

\begin{lstlisting}[language=bash, caption={Service/route registration pattern used for every microservice in \texttt{setup-kongGateway.sh}}]
register_service() {
    local service_name=$1 service_url=$2 route_path=$3 route_name=$4

    curl -s -X POST "http://$addressKong:8001/services/" \
        --data "name=$service_name" --data "url=$service_url" \
        -H "Content-Type: application/x-www-form-urlencoded"

    curl -s -X POST "http://$addressKong:8001/services/$service_name/routes" \
        --data "name=$route_name" --data "paths=$route_path" \
        --data "strip_path=true" \
        -H "Content-Type: application/x-www-form-urlencoded"
}

# e.g.
register_service "prosumer-service" "http://$addressProsumer:8080" "/prosumer" "prosumer-route"
register_service "assetlink-service" "http://$addressAssetLink:8080" "/assetlink" "assetlink-route"
# ... one call per microservice, see Table 1
\end{lstlisting}

\noindent
The \texttt{strip\_path=true} option (added during this sprint's Kong setup
work, see commit \texttt{fix(kong): include strip\_path option}) is what makes
\texttt{/prosumer/Prosumer} resolve to \texttt{/Prosumer} on the upstream
service --- without it, Kong would forward the full original path and every
call would 404 against the microservice's JAX-RS routes.

\subsection{Konga administration UI}

\texttt{KongaTerraform/EC2Konga.tf} follows the same \texttt{t3.small} /
all-open-security-group / \texttt{user\_data} pattern and starts the Konga
container, which was used throughout development to visually confirm that
\texttt{setup-kongGateway.sh} had registered the expected eight services and
eight routes (Table~\ref{tab:kong-routes}) and to inspect/debug routing while
the BPMN service tasks were being wired up.

\subsection{BPMN process definitions --- parametrization conventions}

Two parametrization conventions run through every BPMN file and are the most
important thing to know when redeploying the orchestration layer against a
\emph{new} environment (a fresh \texttt{terraform apply} always produces new
EC2 public DNS names):

\begin{enumerate}
    \item \textbf{\texttt{camundaBaseUrl} process variable} --- every Tasklist
          form-generation service task needs to know the base URL of the
          Camunda REST API it is running against. It must be supplied when
          \emph{starting} a none-start-event process
          (\texttt{"variables": \{"camundaBaseUrl": "http://<camunda-dns>:8080"\}})
          and explicitly \emph{forwarded} inside the \texttt{variables} object
          of every message that starts a collaborating process (Defect~1,
          Section~\ref{sec:bugs}).
    \item \textbf{Kong base URL in HTTP Connector service-task targets} ---
          every outbound REST call embedded in a service task's URL must
          point at \texttt{http://<kong-dns>:8000/<route-path>/...}
          (Table~\ref{tab:kong-routes}), \emph{not} at an individual
          microservice's EC2 address. After a redeploy that produces a new
          Kong DNS, all such URLs across the seven \texttt{.bpmn} files must
          be updated --- which is exactly what the global \texttt{sed}
          substitution in Listing~\ref{lst:sed-fix} automates.
\end{enumerate}

\noindent
Both conventions, the current live endpoint values, and the redeploy command
(Listing~\ref{lst:redeploy}) are kept up to date in
\texttt{BPMN/EndToEnd\_Tests.md}, which doubles as the
parametrization runbook for whoever re-runs the demo after a fresh deployment.

\subsection{Recap: Sprint 1 source code, Terraform, installation and parametrization}

The microservices, shared RDS database, Kafka broker, Docker/Terraform
deployment automation and their parametrization were documented in full in
\texttt{Report.md} (Sections~3--6) and are summarised here for traceability;
they are the components that the BPMN flows in Section~\ref{sec:architecture}
orchestrate through Kong.

\begin{table}[h!]
\centering
\small
\caption{Sprint 1 components orchestrated by the Sprint 2 BPMN flows (full documentation in \texttt{Report.md})}
\label{tab:sprint1-recap}
\begin{tabular}{@{}lll@{}}
\toprule
\textbf{Component} & \textbf{Stack} & \textbf{Documented in \texttt{Report.md}} \\
\midrule
7 domain microservices + AI MS & Quarkus 3.27.2 / Java 17, reactive MySQL & \S2.1, \S3 \\
Shared database & AWS RDS MySQL \texttt{db.t4g.micro}, schema \texttt{VPPaaS} & \S2.3, \S4.1 \\
Event backbone & Apache Kafka 4.1.1 (KRaft, single broker) & \S2.2, \S4.2 \\
AI forecasting runtime & Ollama (\texttt{llama3.2}) on \texttt{t3.large} EC2 & \S2.5, \S4.3 \\
Image registry & Docker Hub & \S2.4 \\
Deployment automation & \texttt{DeploymentAutomation-\{ubuntu,macOS\}.sh} & \S5.3 \\
Parametrization & \texttt{access.sh}, \texttt{application.properties} template & \S6.1, \S6.2, \S6.3 \\
\bottomrule
\end{tabular}
\end{table}

\noindent
\textbf{Installation procedure (end to end, Sprint~1 + Sprint~2):}
\begin{enumerate}
    \item Configure \texttt{access.sh} with AWS credentials, Docker Hub
          credentials and \texttt{JAVA\_HOME}, then \texttt{source access.sh}
          (\texttt{Report.md}~\S5.2).
    \item Run \texttt{./DeploymentAutomation-ubuntu.sh} (or \texttt{-macOS.sh}):
          provisions RDS, Kafka, the seven microservices plus the AI service,
          and --- as extended for this sprint --- the Camunda, Kong and Konga
          EC2 instances, capturing each component's public DNS via
          \texttt{terraform output} and injecting it into the next step's
          configuration (\texttt{Report.md}~\S5.3).
    \item Create the static Kafka topics (\texttt{Report.md}~\S5.4).
    \item Run \texttt{setup-kongGateway.sh} (automatically invoked by the
          deployment script, or manually with the captured DNS values) to
          register the eight Kong services/routes of Table~\ref{tab:kong-routes}.
    \item Deploy the BPMN process definitions to Camunda
          (Listing~\ref{lst:redeploy}, once per \texttt{.bpmn} file) and
          confirm in Operate that all eight \texttt{processDefinitionId}s of
          Table~\ref{tab:flows} are listed under \emph{Processes}.
    \item Run the end-to-end tests of Section~\ref{sec:testing} (or follow
          \texttt{BPMN/EndToEnd\_Tests.md} step by step) to validate
          the full chain \emph{Camunda~$\to$~Kong~$\to$~microservice~$\to$~RDS/Kafka}.
\end{enumerate}

% ===========================================================================
\section{Conclusion}
% ===========================================================================

Sprint~2 closes the gap between the data plane delivered in Sprint~1 and the
business-process layer required by the VPPaaS brief. All eight required
business processes --- \emph{Prosumer Management}, \emph{Utility Operator
Management}, \emph{Asset Link Management}, \emph{Telemetry Ingestion},
\emph{Flexibility Emission}, \emph{Grid Balancing Recommendation},
\emph{Energy Analytics} and \emph{Flexibility Forecasting (AI)} --- were
modelled in BPMN, deployed to Camunda~8, and wired to the underlying
microservices exclusively through the Kong API Gateway, following the
single-entry-point integration pattern the brief requires.

Critically, the flows were not only modelled but \textbf{exercised end to
end} against the live AWS deployment: the \emph{Prosumer Management} flow was
run from a cold Camunda REST API call all the way to a persisted row in the
shared RDS database (Section~\ref{sec:e2e-prosumer}), and the remaining seven
flows were documented with reproducible start commands, Tasklist step
sequences and Kong-side verification calls
(\texttt{BPMN/EndToEnd\_Tests.md}). That process surfaced three
classes of latent defects --- a missing cross-process variable, a systemic
placeholder left unresolved in 21 service-task URLs across all seven BPMN
files, and a variable-name case mismatch silently dropping a business field
--- all of which were corrected at the BPMN-definition level and re-validated
with fresh process instances (Section~\ref{sec:bugs}). Finding and fixing
these issues is, in itself, the strongest evidence that the integration layer
was genuinely tested rather than merely modelled: none of the three defects
was visible from a static reading of the diagrams, and all three would have
caused the live demonstration to fail had they not been caught beforehand.

Finally, this report documents --- alongside the new orchestration-layer
artefacts (BPMN models, Camunda/Kong/Konga Terraform, Kong service-registration
script, parametrization conventions for \texttt{camundaBaseUrl} and the Kong
base URL) --- the full installation procedure that takes the platform from
zero to a running, end-to-end-testable system, building on and extending the
Sprint~1 documentation already produced in \texttt{Report.md}.

\end{document}
